% GENERATED FILE. Edit canonical lessons, then run npm run generate:books.
% canonical-source-hash: fnv1a64:22e8db5bbdd0e5c7
% canonical-lessons: GE-C478-reue, GE-C478-bereuen, GE-C478-verzeihen, GE-C478-beleidigen, GE-C478-kraenken

\chapter{Reue und Verzeihung}
\label{ch:ge-reue-und-verzeihung}
\hlchaptermodality{478}

\begin{chapteropening}
\textbf{By the end of this chapter:} \emph{I can talk about regret, hurt feelings and forgiving.}

Complete the last lesson of chapter 478: I can talk about regret, hurt feelings and forgiving.
\end{chapteropening}

\section[die Reue]{die Reue — remorse, regret}

\label{lesson:GE-C478-reue}



\begin{quote}
Before the new one: say the German for reckless, then the German for to dare.
\end{quote}

\subsection*{You'll want to know: die Reue}
\textbf{die Reue} — \textquotedblleft{}remorse, regret\textquotedblright{}.

\begin{quote}
Der Täter zeigte keine Reue.
\end{quote}

\begin{quote}
\emph{The culprit showed no remorse.}
\end{quote}

No plural.

\begin{cousinweb}
\emph{Reue} is inherited from Old High German \emph{riuwa}, from Proto-Germanic \emph{hrewwō} \textquotedblleft{}pain, sadness, regret\textquotedblright{} (EtymDB); English \emph{rue} is a cousin.
\end{cousinweb}

\begin{grammarlens}[title={What you've built}]
One more word for telling, explaining and describing at B1.
\end{grammarlens}

\subsection*{Your turn}
\begin{itemize}
\raggedright
  \item \emph{Say it:} \emph{die Reue}
  \item \emph{Say it:} \emph{die Reue}, once more
  \item \emph{Say it:} a sentence of your own with \emph{die Reue}
  \item \emph{Recall:} say the German for reckless, then the German for to dare, then say \emph{die Reue} again
\end{itemize}

\begin{tcolorbox}[breakable,colback=teal!4,colframe=teal!35!black,title={Before you move on}]
What does die Reue mean? (\textquotedblleft{}Remorse, regret\textquotedblright{}.) Say it once more.
\end{tcolorbox}
\section[bereuen]{bereuen — to regret}

\label{lesson:GE-C478-bereuen}



\begin{quote}
Before the new one: say the German for to dare, then the German for remorse.
\end{quote}

\subsection*{You'll want to know: bereuen}
\textbf{bereuen} — \textquotedblleft{}to regret\textquotedblright{}.

\begin{quote}
Sie bereut, dass sie so wenig gelernt hat.
\end{quote}

\begin{quote}
\emph{She regrets studying so little.}
\end{quote}

Inseparable: \emph{hat bereut}.

\begin{cousinweb}
\emph{bereuen} is built on \emph{be-} and the old verb \emph{reuen} \textquotedblleft{}to cause sorrow\textquotedblright{}, from the same source as \emph{Reue}, Proto-Germanic \emph{hrewwō} \textquotedblleft{}pain, regret\textquotedblright{} (EtymDB).
\end{cousinweb}

\begin{grammarlens}[title={What you've built}]
One more word for telling, explaining and describing at B1.
\end{grammarlens}

\subsection*{Your turn}
\begin{itemize}
\raggedright
  \item \emph{Say it:} \emph{bereuen}
  \item \emph{Say it:} \emph{bereuen}, once more
  \item \emph{Say it:} a sentence of your own with \emph{bereuen}
  \item \emph{Recall:} say the German for to dare, then the German for remorse, then say \emph{bereuen} again
\end{itemize}

\begin{tcolorbox}[breakable,colback=teal!4,colframe=teal!35!black,title={Before you move on}]
What does bereuen mean? (\textquotedblleft{}To regret\textquotedblright{}.) Say it once more.
\end{tcolorbox}
\section[verzeihen]{verzeihen — to forgive}

\label{lesson:GE-C478-verzeihen}



\begin{quote}
Before the new one: say the German for remorse, then the German for to regret.
\end{quote}

\subsection*{You'll want to know: verzeihen}
\textbf{verzeihen} — \textquotedblleft{}to forgive\textquotedblright{}.

\begin{quote}
Sie hat ihm den Fehler sofort verziehen.
\end{quote}

\begin{quote}
\emph{She forgave him the mistake at once.}
\end{quote}

Irregular: \emph{verzieh, hat verziehen}. \emph{Verzeihung!} — Excuse me!

\begin{cousinweb}
\emph{verzeihen} comes from Middle High German \emph{verzīhen} (EtymDB), built on \emph{zīhen} \textquotedblleft{}to accuse\textquotedblright{}, from Proto-Germanic \emph{tīhaną} (EtymDB); it is usually explained as first meaning \textquotedblleft{}to give up a claim\textquotedblright{}.
\end{cousinweb}

\begin{grammarlens}[title={What you've built}]
One more word for telling, explaining and describing at B1.
\end{grammarlens}

\subsection*{Your turn}
\begin{itemize}
\raggedright
  \item \emph{Say it:} \emph{verzeihen}
  \item \emph{Say it:} \emph{verzeihen}, once more
  \item \emph{Say it:} a sentence of your own with \emph{verzeihen}
  \item \emph{Recall:} say the German for remorse, then the German for to regret, then say \emph{verzeihen} again
\end{itemize}

\begin{tcolorbox}[breakable,colback=teal!4,colframe=teal!35!black,title={Before you move on}]
What does verzeihen mean? (\textquotedblleft{}To forgive\textquotedblright{}.) Say it once more.
\end{tcolorbox}
\section[beleidigen]{beleidigen — to insult, to offend}

\label{lesson:GE-C478-beleidigen}



\begin{quote}
Before the new one: say the German for to regret, then the German for to forgive.
\end{quote}

\subsection*{You'll want to know: beleidigen}
\textbf{beleidigen} — \textquotedblleft{}to insult, to offend\textquotedblright{}.

\begin{quote}
Mit seinem Witz hat er die Gastgeberin beleidigt.
\end{quote}

\begin{quote}
\emph{He offended the hostess with his joke.}
\end{quote}

\emph{Beleidigt sein} — to be offended, to sulk. The noun: \emph{die Beleidigung}.

\begin{cousinweb}
\emph{beleidigen} is built on \emph{be-} and \emph{Leid} \textquotedblleft{}harm, sorrow\textquotedblright{}, which is inherited from Old High German \emph{leid}, from Proto-Germanic \emph{laiþą} (EtymDB).
\end{cousinweb}

\begin{grammarlens}[title={What you've built}]
One more word for telling, explaining and describing at B1.
\end{grammarlens}

\subsection*{Your turn}
\begin{itemize}
\raggedright
  \item \emph{Say it:} \emph{beleidigen}
  \item \emph{Say it:} \emph{beleidigen}, once more
  \item \emph{Say it:} a sentence of your own with \emph{beleidigen}
  \item \emph{Recall:} say the German for to regret, then the German for to forgive, then say \emph{beleidigen} again
\end{itemize}

\begin{tcolorbox}[breakable,colback=teal!4,colframe=teal!35!black,title={Before you move on}]
What does beleidigen mean? (\textquotedblleft{}To insult, to offend\textquotedblright{}.) Say it once more.
\end{tcolorbox}
\section[kränken]{kränken — to hurt (someone's feelings)}

\label{lesson:GE-C478-kraenken}



\begin{quote}
Before the new one: say the German for to forgive, then the German for to insult.
\end{quote}

\subsection*{You'll want to know: kränken}
\textbf{kränken} — \textquotedblleft{}to hurt (someone's feelings)\textquotedblright{}.

\begin{quote}
Seine harten Worte haben sie tief gekränkt.
\end{quote}

\begin{quote}
\emph{His harsh words hurt her deeply.}
\end{quote}

Of feelings only; a hurt body is \emph{verletzt}.

\begin{cousinweb}
\emph{kränken} is built on \emph{krank} \textquotedblleft{}sick\textquotedblright{}, which is inherited from Old High German \emph{krank}, from Proto-Germanic \emph{krangaz} (EtymDB): to make someone feel ill.
\end{cousinweb}

\begin{grammarlens}[title={What you've built}]
One more word for telling, explaining and describing at B1.
\end{grammarlens}

\subsection*{Your turn}
\begin{itemize}
\raggedright
  \item \emph{Say it:} \emph{kränken}
  \item \emph{Say it:} \emph{kränken}, once more
  \item \emph{Say it:} a sentence of your own with \emph{kränken}
  \item \emph{Recall:} say the German for to forgive, then the German for to insult, then say \emph{kränken} again
\end{itemize}

\begin{tcolorbox}[breakable,colback=teal!4,colframe=teal!35!black,title={Before you move on}]
What does kränken mean? (\textquotedblleft{}To hurt (someone's feelings)\textquotedblright{}.) Say it once more.
\end{tcolorbox}
